\section{Szczegółowa specyfikacja projektu i opis wymagań}

\subsection{Cel i kryterium sukcesu}

Sukces projektu definiowany jest jako uzyskanie klasyfikatora (lub metody statystycznej) osiągającego \textbf{dokładność co najmniej $51\%$} na wcześniej odłożonym zbiorze testowym, wraz z opisem warunków, w których wynik jest miarodajny (długość tekstu, język, źródło generatora).

Dodatkowym kryterium jakości jest udokumentowanie przypadków, w których metoda zawodzić będzie przewidywalnie (teksty krótkie, silna parafraza, mieszanka edycji ludzkiej i AI).

\subsection{Wejście / wyjście}

\begin{itemize}
	\item \textbf{Wejście:} tekst w formie ciągłej (plik lub łańcuch znaków) spełniający minimalną długość.
	\item \textbf{Wyjście:} etykieta \{\texttt{human}, \texttt{ai}\} oraz - jeśli to możliwe - wynik liczbowy (score / prawdopodobieństwo).
	\item \textbf{Metadane pomocnicze (opcjonalnie):} długość w tokenach/słowach, język, nazwa użytej metody.
\end{itemize}

\subsection{Przepływ przetwarzania}

\begin{enumerate}
	\item Wczytanie tekstu i normalizacja (usunięcie zbędnych białych znaków, ewentualna detekcja języka).
	\item Sprawdzenie minimalnej długości; poniżej progu zwracany jest komunikat o braku miarodajności zamiast twardej etykiety.
	\item Ekstrakcja cech (statystycznych) albo wnioskowanie modelu ML.
	\item Podjęcie decyzji klasyfikacyjnej według ustalonego progu.
	\item Zapis wyniku i podstawowych metryk przy ewaluacji wsadowej.
\end{enumerate}

\subsection{Założenia}

\begin{itemize}
	\item Prace rozpoczynają się od podejścia statystycznego (zgodnie z konsultacją eksperta przy małej skali danych).
	\item Dla wariantu ML teksty treningowe powinny być możliwie zbliżone tematycznie (rekomendacja eksperta).
	\item Limit długości tekstu zostanie ustalony empirycznie.
	\item Projekt nie korzysta z prywatnego detektora znaków wodnych OpenAI; analizowana jest detekcja niezależna od współpracy generatora.
	\item Zbiór testowy nie jest używany do strojenia progów ani hiperparametrów.
\end{itemize}

\subsection{Ograniczenia}

\begin{itemize}
	\item Dokładność $51\%$ jest progiem minimalnym akceptacji, nie celem optymalnym.
	\item Wynik może nie przenosić się na inne domeny, modele LLM ani język polski, jeśli trening odbywał się głównie na danych angielskich.
	\item Metoda post-hoc nie zastępuje watermarkingu dostawcy i nie identyfikuje konkretnego modelu ani użytkownika.
	\item Brak odporności na celową parafrazę lub ciężką redakcję ludzką jest akceptowany na tym etapie.
\end{itemize}

\subsection{Plan ewaluacji}

\begin{itemize}
	\item Podział danych: trening / walidacja / test (lub trening / test przy małym zbiorze).
	\item Metryki: dokładność; dodatkowo precyzja, czułość lub F1 dla klasy \texttt{ai}.
	\item Analiza wrażliwości na długość tekstu (kilka progów długości).
	\item Krótki opis błędów: przykłady fałszywie dodatnich i fałszywie ujemnych.
\end{itemize}

\subsection{Środowisko realizacji}

\begin{itemize}
	\item Język implementacji: Python.
	\item Biblioteki: typowy stos analizy tekstu / ML (np.\ scikit-learn, transformers - zależnie od wariantu).
	\item Sprzęt: komputer studencki; ewentualnie darmowe środowisko notebookowe z GPU.
\end{itemize}
